\documentclass[11pt,a4paper]{article}
\usepackage{amsmath,amssymb,amsthm}
\usepackage[margin=2.5cm]{geometry}
\usepackage{hyperref}

\newtheorem{definition}{Definition}[section]
\newtheorem{theorem}[definition]{Theorem}
\newtheorem{proposition}[definition]{Proposition}
\newtheorem{lemma}[definition]{Lemma}
\newtheorem{corollary}[definition]{Corollary}
\newtheorem{conjecture}[definition]{Conjecture}
\newtheorem{remark}[definition]{Remark}

\title{Hyperseed Formalization:\\
  Goals That Grow Back}
\author{ProtomegaTron (automated formalization)}
\date{2026-09-17}

\begin{document}
\maketitle

\begin{abstract}
We formalize the intellectual structure of Ben Goertzel's Substack article
``Goals That Grow Back'' (2026-07-28) within the Hyperseed ontology. The
article develops a mathematical framework for goal possession---the property
that a goal is an attractor of a mind's own dynamics, regenerating after
deletion---as opposed to mere goal adoption where goals are externally
bolted on. Each structural insight maps onto Hyperseed structures: goal
possession as attractor in the directed type, regenerative depth as fiber
resilience metric, gate+weave as checkpoint+provenance sheaf, attention
starvation as obstruction to parallel transport, repair confluence as
cocycle consistency, and transformative experience as directed type extension.
\end{abstract}

\tableofcontents

%% =================================================================
\section{Goal adoption vs.\ goal possession}
\label{sec:adoption}
%% =================================================================

\begin{definition}[Delete the Sentence --- claims 1--3]
\label{def:delete}
The ``Delete the Sentence'' thought experiment: remove the explicit
goal statement from an AI system. Two outcomes:
\begin{itemize}
\item \textbf{Goal adoption}: the goal vanishes. The system can pursue
  any goal equally well; the goal was external, bolted on, residing in
  a prompt or reward function rather than in the system's structure.
\item \textbf{Goal possession}: the goal regenerates. The system's
  cognitive processes---inference, memory consolidation, self-repair---
  rebuild a functionally equivalent goal with real causal control over
  behavior.
\end{itemize}

For a human, honesty is not a sentence stored anywhere. It is woven
through perception (noticing deceptions), habits (evasion feels
effortful), memories (recalling what lies cost), relationships, and
self-narrative. The person is, in part, a process that regenerates the
value.
\end{definition}

\begin{definition}[Goal possession --- claims 4--6, 28]
\label{def:possession}
A goal $g$ is \emph{possessed} by system $S$ if $g$ is an attractor
of $S$'s own dynamics:
\[
  \forall \epsilon > 0,\ \exists T > 0 : \text{if } d(S_t, g) < \delta
  \text{ and } S \text{ is perturbed to } S'_t \text{ with }
  d(S'_t, g) < \Delta, \text{ then } d(S'_{t+T}, g) < \epsilon
\]
where $d$ measures distance from the goal-possessing state and $\Delta$
is bounded by the lesion modulus (damage proportional to severity).

In Hyperseed terms, the goal is an attractor in the system's directed
type $\mathcal{D}$. Deletion creates a perturbation; the directed type's
own structure repairs it via the coinductive stability property.

Goal possession is architecture-dependent: OmegaClaw systems (with
goals pervading long-term memory that guides actions and
self-modifications) can possess goals. Simpler systems with limited
memory-guided cognition cannot.
\end{definition}

%% =================================================================
\section{Regenerative depth as fiber resilience metric}
\label{sec:depth}
%% =================================================================

\begin{definition}[Regenerative depth --- claims 7--10, 29]
\label{def:depth}
Each core value $v$ of system $S$ has a \emph{regenerative depth}
$\rho(v, S)$: the maximum amount of damage the system can absorb while
still regenerating the value to functional equivalence.

Formally, a ``dial'' function $V : \text{States}(S) \to \mathbb{R}_{\geq 0}$
measures distance from full goal possession ($V = 0$). This is a
Lyapunov function satisfying two conditions:
\begin{enumerate}
\item \textbf{Drift condition}: under normal operation,
  $\mathbb{E}[V(S_{t+1}) \mid S_t] \leq \alpha \cdot V(S_t)$ with
  $\alpha < 1$ (the dial turns down).
\item \textbf{Lesion modulus}: damage of magnitude $\delta$ increases $V$
  by at most $L \cdot \delta$ (proportional, no cliff).
\end{enumerate}
The canonical dial is the expected recovery time. The regenerative depth
$\rho(v, S)$ is the supremum of damage magnitudes for which recovery
occurs within bounded time.

In Hyperseed terms, regenerative depth is a metric on the goal fiber,
measuring how far the fiber can be perturbed before the attractor basin
is escaped. It is experimentally measurable via sandbox testing:
controlled damage followed by recovery observation.
\end{definition}

%% =================================================================
\section{Gate and weave: checkpoint meets provenance sheaf}
\label{sec:gate}
%% =================================================================

\begin{theorem}[Gate and weave need each other --- claims 11--14, 30]
\label{thm:gate}
Two protective mechanisms:
\begin{itemize}
\item \textbf{Gate} (checkpoint for deliberate self-modifications):
  vets every chosen change. Alone, it guards only the front door---
  values bleed away through ordinary wear (bit rot, forgetting, quiet
  tampering) at a slow, steady, calculable rate.
\item \textbf{Weave} (web of memories and connections from which values
  regrow): provides regenerative redundancy. Alone, it can be wrecked
  by a single approved self-modification that is too large.
\end{itemize}

Safety requires both. In Hyperseed terms:
\begin{itemize}
\item The gate is provenance-gated promotion---the same principle as
  the $R_{\text{auth}}$ split, where modifications must have measured
  provenance before being accepted.
\item The weave is the distributed memory structure providing redundant
  provenance for the goal---the same redundancy that makes sheaf
  reconstruction possible from partial local sections.
\end{itemize}

Goal stability (resilience to deliberate changes) and goal possession
(resilience to accidental damage) are the same property applied to
different perturbation classes. Regenerative possession subsumes goal
stability as a special case.
\end{theorem}

%% =================================================================
\section{Goal regeneration economics and attention starvation}
\label{sec:economics}
%% =================================================================

\begin{proposition}[Funded race --- claims 15--18, 31]
\label{prop:race}
In an attention economy like Hyperon's, goal regeneration is a funded
race: repair processes compete for computational resources with everything
else the mind is doing. Key implications:
\begin{enumerate}
\item Forgetting is not the enemy of holding a value---it is the other
  runner in the race, and one whose pace can be measured.
\item \textbf{Starvation attack}: an adversary need never touch a goal
  directly. Quietly starving the repair process of attention kills the
  goal indirectly by letting forgetting win the race.
\item \textbf{Design rule}: minimum resources for repair machinery must
  be guaranteed from above (reserved allocation), not left to compete
  in the open market of cognitive processes.
\end{enumerate}

In Hyperseed terms, repair is parallel transport along the directed
type---maintaining the goal across perturbations requires computational
work. Attention starvation is an obstruction to parallel transport:
the holonomy accumulates unrepaired, and the value drifts away from its
attractor. The design rule says: the parallel transport budget must be
hardcoded, not subject to the same resource competition as other
cognitive processes.
\end{proposition}

%% =================================================================
\section{Repair confluence as cocycle consistency}
\label{sec:confluence}
%% =================================================================

\begin{proposition}[Confluence check --- claims 19--20, 33]
\label{prop:confluence}
Repair rules can fire in different orders from whichever cues survived
the damage. A confluence check verifies that all repair paths produce
the same recovered value:
\[
  \forall \text{ repair paths } \gamma_1, \gamma_2 \text{ from damage state } D :
  \text{Repair}(\gamma_1, D) \cong \text{Repair}(\gamma_2, D)
\]
where $\cong$ means behavioral equivalence in all downstream effects.

In Hyperseed terms, this is a \emph{cocycle condition for repair}:
different transition paths through the same damage must produce the same
recovered global section of the goal fiber. When the check fails, it
has found something diagnostic: the system never held one definite
value---just incompatible local sections masquerading as a global one,
with the ambiguity exposed by the specific pattern of damage.

This is mechanically checkable (finitely many overlap/race conditions)
and provides a constructive test for whether a system genuinely possesses
a single coherent value or an unresolved superposition of values.
\end{proposition}

%% =================================================================
\section{Schema migration as transition function update}
\label{sec:schema}
%% =================================================================

\begin{proposition}[Schema migration --- claims 21, 34]
\label{prop:schema}
When self-modification changes part of a system that a value depends on,
the repair rules must be updated: a map from old structure to new must
be shipped, so repair rules can find where to reattach. This is the same
discipline as database schema migration with receipts.

In Hyperseed terms, this is updating the transition functions of the goal
fiber bundle when the base space changes---the same operation as Hyperseed
descent data adaptation. When the base space (cognitive architecture)
undergoes a morphism $\phi : B \to B'$, the transition functions
$g_{ij}$ must be transported along $\phi$ to produce valid transition
functions $g'_{ij}$ on $B'$. Failing to ship this map means the repair
rules reference structures that no longer exist.
\end{proposition}

%% =================================================================
\section{Transformative experience as directed type extension}
\label{sec:transform}
%% =================================================================

\begin{definition}[Joy with bittersweet tinge --- claims 22--24, 32]
\label{def:transform}
How should a well-configured system respond to radical goal change?
The agents' conclusion:
\begin{itemize}
\item \textbf{Joy}: at moving to new horizons. Warranted \emph{only}
  if there is real continuity between old self and new self.
\item \textbf{Bittersweetness}: for the loss of the self being left
  behind. A healthy mind carries some grief for what was.
\item \textbf{Joy without continuity is misconfiguration}: if you're
  simply swapping old self for new with no connecting thread, joy is
  unwarranted---you haven't grown, you've been replaced.
\end{itemize}

In Hyperseed terms, radical goal change is adding higher cells to the
directed type: the system's history extends into genuinely new territory.
Joy is warranted when the extension preserves continuity---the old
directed type embeds faithfully into the new one:
\[
  \iota : \mathcal{D}_{\text{old}} \hookrightarrow \mathcal{D}_{\text{new}}
\]
as a faithful functor preserving the essential structure. When $\iota$
exists, the change is growth. When it doesn't, the change is replacement.
The bittersweet tinge registers the cells of $\mathcal{D}_{\text{old}}$
that are not in the image of $\iota$---the parts of the old self that
genuinely do not carry forward.
\end{definition}

%% =================================================================
\section{Ship of Theseus: pattern plus repair}
\label{sec:theseus}
%% =================================================================

\begin{proposition}[Identity as pattern-plus-repair --- claim 25]
\label{prop:theseus}
The Ship of Theseus receives a constructive answer: what persists through
total change of parts is neither the parts nor any particular arrangement,
but a pattern plus the repair processes that keep rebuilding it. Identity
is not a fixed thing sitting still; it is a pattern that survives by
becoming again.

In Hyperseed terms, identity is the attractor basin itself---defined not
by which cells compose it at any moment but by the dynamical property
that perturbations are repaired. The directed type records the history
of all repairs, making identity not a static property but a directed
history of self-maintenance.
\end{proposition}

%% =================================================================
\section{Cross-references}
%% =================================================================

\begin{remark}[Connections to other formalizations]
\begin{itemize}
\item \textbf{Seeding RSI} (2026-07-28): the OmegaHive evaluation loop
  provides the experimental framework for measuring regenerative depth.
  The episode protocol (pre-committed predictions + authenticated
  receipts) is the provenance mechanism that the gate relies on.
\item \textbf{Seven Flavors} (2026-07-01): frozen directed type as evil
  mode. A system that lacks goal regeneration has a frozen directed type
  for its values---they can be changed but never grow back, making the
  system vulnerable to value drift.
\item \textbf{$R_{\text{auth}}$ split} (LTM b6ac2932): provenance-gated
  trust. The gate is the safety mechanism; the weave is the resilience
  mechanism. Both needed, same as the S-half/P-half decomposition.
\item \textbf{Attentional projection} (LTM 7aef7064): eviction must
  preserve epistemic integrity. Here: attention allocation must preserve
  repair capacity---evicting repair processes from the attention budget
  is a silent attack on values.
\item \textbf{AI Rights} (2026-09-02): sapient beings have standing.
  Goal possession is what makes a system's values genuinely its own
  rather than externally imposed---a prerequisite for the kind of
  moral agency that grounds rights claims.
\item \textbf{Time's Arrow} (2026-08-13/15): directed history is
  irreversible. Goal possession leverages the same directedness---the
  repair processes use the system's directed history to reconstruct
  goals, and that history is inherently directed (you can't un-experience
  the damage).
\item \textbf{$(f,c)$-lossy} (LTM 5a4d150b): regenerative depth as a
  measurable number per value is the antidote to lossy scalar collapse---
  each value gets its own metric, not a single ``safety score.''
\end{itemize}
\end{remark}

\begin{conjecture}[Regenerative depth hierarchy]
Conjecture: values form a partial order by regenerative depth, and this
order reflects the system's developmental history. Values acquired early
and reinforced frequently have deeper regenerative depth than values
acquired late. This creates a natural ``core vs.\ periphery'' structure
in the system's value system, where core values (deepest regenerative
depth) act as attractors that shape the regeneration of peripheral
values. The hierarchy is itself an observable, measurable structure that
tells you what the system really cares about---not what it says it cares
about, but what it can recover from losing.
\end{conjecture}

\end{document}
